\chapter{Límites y continuidad}

\section{Límites de sucesiones}

\begin{definition}\label{def:limit}
Una sucesión $(a_n)$ converge a $L\in\mathbb{R}$ si para todo $\varepsilon>0$
existe $N\in\mathbb{N}$ tal que $|a_n-L|<\varepsilon$ para todo $n\ge N$.
\end{definition}

\begin{theorem}[Unicidad del límite]\label{thm:uniqueness}
Si $(a_n)$ converge, su límite es único.
\end{theorem}
\begin{proof}
Si $L\neq L'$ fueran límites, tomando $\varepsilon=|L-L'|/2$ en la
definición~\ref{def:limit} se llega a una contradicción.
\end{proof}

\begin{example}
La sucesión $a_n = 1/n$ converge a $0$.
\end{example}

\section{Funciones continuas}

\begin{lemma}\label{lem:comp}
La composición de funciones continuas es continua.
\end{lemma}

\begin{figure}[ht]
\centering
\begin{tikzpicture}
\begin{axis}[width=7cm,height=5cm,xlabel={$x$},ylabel={$f(x)$},domain=-2:2,samples=60]
\addplot[thick,rm-series-1]{x^2};
\addplot[thick,dashed,rm-guide,domain=-1:2]{2*x-1};
\addplot[only marks,mark=*,rm-highlight] coordinates {(1,1)};
\end{axis}
\end{tikzpicture}
\caption{La función $f(x)=x^2$.}\label{fig:parabola}
\end{figure}

Véase la figura~\ref{fig:parabola} y el teorema~\ref{thm:uniqueness}; también \cite{rudin,cauchy}.
